\section{HAL-MLE for Pathwise Differentiable Statistical Estimands} \label{section:HAL-MLE for Pathwise Differentiable Statistical Estimands}
One could use the density estimator for statistical estimands such as moments, the median, the CDF, or the survival probability. Let $\Psi:{\cal M}\rightarrow \R$ be a pathwise differentiable target parameter with canonical gradient (also called efficient influence curve) $D^*_P$ at $P$ where ${\cal M}=\{P_f: f\in D^{(k)}_U([0,1]^d \}$, then, for a path $P_{\varepsilon}^h$ through $P$ with direction $h$, $\frac{\mathrm{d}}{\mathrm{d} \varepsilon} \Psi(P_{\varepsilon}^h) = \mathbb{E}_P[ D^*_P \, h]$. Let $R(P, P_0)\equiv \Psi(P)-\Psi(P_0)+P_0D^*_P$ be the exact remainder implied by the canonical gradient, which is known to be a second-order difference due to it being the exact remainder in a first order Tailor expansion $R(P,P_0)=\Psi(P)-\Psi(P_0)-(P-P_0)D^*_P$.
We have $\Psi(P)=\Psi(P_f)=\Psi^F(f)$ for a $\Psi^F: D^{(k)}_U([0,1]^d)\rightarrow \R$.
We can represent $D^*_P=D^*_f$ with $f=f(P)$.

\subsection{"Naive" Approach: Plug-in} \label{Plug-in HAL-MLE}
Plugging in the density would be a natural approach, and we will show that it yields an asymptotically efficient estimator, possibly achieved by using undersmoothing. We here provide a standard TMLE-type analysis establishing asymptotic efficiency under the assumption that $D^{(k)}({\cal R}_n)$ approximates $D^{(k)}_U([0,1])$. 

The tangent space generated by paths through $f$ is given by closure of linear span $\{S_f(\phi): \phi\in D^{(k)}_U([0,1])\}$. Thus, we can represent 
$D^*_f=S_f(\phi_f)$ for some $\phi_f\in D^{(k)}_M([0,1])$ for $\forall k$ including 0. Therefore, if the basis ${\cal R}_n$ is chosen so that $D^{(k)}({\cal R}_n)$  yields a $O^+(1/J_n^{k+1})$ $L^2$ approximation of $D^{(k)}_U([0,1])$, then we can find a projection $\tilde{\phi}_f$ of $\phi_f$ onto $D^{(k)}({\cal R}_n)$ that approximates $\phi_f$ in $L^2$-norm within distance $O^+(1/J_n^{k+1})$. 


For simplicity, firstly, consider the case that $f_n$ is a relaxed HAL-MLE so that it equals an MLE over $D^{(k)}({\cal R}_n)$, while we assume that it has been established that $\Vert \beta_n\Vert=o_p(1)$ as it would be for the HAL-MLE. Then, $P_n S_{f_n}(\phi_j)=0$ for all $j\in {\cal R}_n$. Let $\tilde{D}_{f_n}$ be the projection in $L^2(\mu)$ of $D^*_{f_n}$ onto $\{S_{f_n}(\phi_j):j\in {\cal R}_n\}$. Then, $P_n \tilde{D}_{f_n}=0$. In that case if $\Vert \tilde{D}_{f_n}-D^*_{f_n}\Vert_{\mu}=o_p^+(1/J_n^{k+1})$, then it follows that $P_n D^*_{f_n}=o_p(n^{-1/2})$ and the efficiency theorem below  applies. For the HAL-MLE we already know we control $\Vert \beta_n\Vert_1=O(1)$, but we would need to show that the scores $P_n S_{f_n}(\phi_j)$, $j\in {\cal R}_n$ are solved good enough so that $P_n \tilde{D}_{f_n}=o_p(n^{-1/2})$. In Section S.6 of the Supplementary Material \citep{hou2026supplement}, we show that this can generally be arranged by extending the approximation condition to the $L_1$ constrained score space.


\begin{theorem}[Efficient Influence Curve Approximation]
\label{thm:Efficient Influence Curve Approximation}
Suppose the sieve \(D^{(k)}(\mathcal R_n)\) satisfies the \(L^2\)‐approximation and it can be extended to its corresponding $L_1$ constrained score space, and maintain the Smoothness Assumptions in Theorem~\ref{thm:asymptotic-linearity-hal-mle}. 
Let \(f_n\) be the cross‐validated HAL–MLE over \(D^{(k)}(\mathcal R_n)\). Then
\[
P_n\,D^*_{f_n}
= o_p\!\bigl(n^{-1/2}\bigr).
\]
\end{theorem}

This directly leads to the asymptotic efficiency theorem. 
\begin{theorem}[Asymptotic Efficiency of Plug‐in HAL‐MLE]
\label{thm:plugin_HAL_asymptotic_efficiency}
Suppose the following conditions hold:
\begin{enumerate}
  \item (L\(^2\)‐Approximation) \(D^{(k)}(\mathcal R_n)\) satisfies the \(L^2\)‐approximation and it can be extended to its corresponding $L_1$ constrained score space.
  \item (Negligible Remainder) \(R(P_{f_n},P_0)=o_p\bigl(n^{-1/2}\bigr)\).
  \item (\(L^2(P_0)\)-continuity of EIC) \(P_0\{(D^*_{f_n}-D^*_{f_0})^2\}
      = o_p\bigl(1\bigr)\).
\end{enumerate}
Then the plug‐in estimator \(\Psi^F(f_n)\) satisfies
\[
\Psi^F(f_n) - \Psi^F(f_0)
= (P_n - P_0)\,D^*_{f_0} \;+\; o_p\!\bigl(n^{-1/2}\bigr),
\]
and is therefore asymptotically efficient.
\end{theorem}

We discussed the plausibility of the first assumption in Section~\ref{subsect: Pointwise Asymptotic Normality of the HAL-MLE}. And here we demonstrate that the second and third assumptions are not only reasonable but also very weak with the examples of survival probability and moments. We leave the argument of median, which is slightly complicated, to Section S.6 of the Supplementary Material \citep{hou2026supplement}. 
For survival probability at $x_0$, we have 
\[P_0{(D^*_{f_n} - D^*_{f_0})} = \int I(x>x_0)(p_{f_n} - p_{f_0}) \mathrm{d}x =O(1)o_p(\Vert p_{f_n} - p_{f_0} \Vert_{\mu}) = o_p(n^{-\frac{k+1}{2k+3}}),
\]
by the Cauchy-Schwartz inequality. This implies \(P_0\{(D^*_{f_n}-D^*_{f_0})^2\} =o_p(n^{-\frac{2k+2}{2k+3}}) = o_p\bigl(1\bigr)\). 
\[R(P_{f_n}, P_0) \equiv \Psi(P_{f_n}) - \Psi(P_0) + P_0 D_P = S_{f_n}(x_0) - S_0(x_0) + P_0 (I(x >x_0) - P_{f_n}(x>x_0)) = 0.
\] 
The argument is basically the same for moments, \[P_0{(D^*_{f_n} - D^*_{f_0})} =\int x^k (p_{f_n} - p_{f_0}) \mathrm{d}x = o_p(n^{-\frac{k+1}{2k+3}}),\] by Cauchy-Schwartz inequality. \[R(P_{f_n}, P_0) = \mathbb{E}_{P_{f_n}}(x^k) - \mathbb{E}_{P_{0}}(x^k) + P_0 (x^k - \mathbb{E}_{P_{f_n}}(x^k)) = 0.\]


\subsection{HAL-TMLE} \label{HAL-TMLE}
We can exploit the information about the target parameter more effectively by incorporating the TMLE framework, namely HAL-TMLE. We found that the targeting step is very compatible with the link function, and each target parameter corresponds to a different targeting. 

\textbf{Targeting Step:} We construct a one-dimensional submodel $\{p_{f_{\varepsilon}^h}:\varepsilon\}$ indexed by a direction $h$ of a path $\{f_{\varepsilon}^h:\varepsilon\}$ through $f$ that satisfies
\begin{align} \label{lfm}
  \left  .  \frac{\mathrm{d}}{\mathrm{d}\varepsilon} \log p_{f_{\varepsilon}^h}\right |_{\varepsilon=0} = D^*_f(x),
\end{align}

\textbf{MLE Step:} Specifically, we select $f_{\varepsilon}^h = f + \varepsilon h$, and thus we have
\[
p_{f_{\varepsilon}^h} = \frac{\exp\bigl(f + \varepsilon h\bigr)}{\displaystyle \int \exp\bigl(f + \varepsilon h\bigr)\mathrm{d}x},
\]
where $f$ plays the role of $f_{n, \beta_n(M)}$. 
We can simply select $h=h_f=D^*_f$ to make this path a local least favorable path (LLFP). Let's denote the resulting path with $f_{n,\varepsilon}$
Then, the first step TMLE update would be defined by 
$f_{n,\varepsilon_n^1}$ with $\varepsilon_n^1=\arg\min_{\varepsilon}-P_n \log p_{f_{n,\varepsilon}}$. One can set $f_n^1=f_{n,\varepsilon_n^1}$ and define $f_{n,\varepsilon}^1$ as this LLFP but now with off-set $f_n^1$, and define $\varepsilon_n^2=\arg\min_{\varepsilon}-P_n \log p_{f_{n,\varepsilon_n^2}^1}$. We can iterate this TMLE-updating till $P_n D^*_{f_n^k}\leq \sigma_n/(n^{1/2}\log n)$, where
$\sigma_n^2 = (P_n D^*_{f_n})^2$ is the estimator of the sample variance of the efficient influence curve at the initial estimator. The final update satisfying this criterion is then the iterative TMLE $f_n^*$, implying the plug-in TMLE $\Psi^F(f_n^*)$ of $\Psi^F(f_0)$. 

A theoretical concern in iterative TMLE is the preservation of the function class.
The HAL-MLE $f_n$ belongs to a function space whose induced distribution $P_{f_n}^0$ lies in a 
Donsker class with covering number equivalent to that of $D^{(k)}([0,1])$. When applying a finite number of targeting steps, this property is preserved. However, performing infinitely many TMLE iterations may drive the estimator outside of this Donsker class.

According to \citet{van2017generally}, a single targeting step already suffices, because the initial HAL-MLE $f_n$ converges to $f_0$ at a rate faster than $n^{-1/4}$ and we also satisfy the positivity condition, the single-step update $f_n^1$ obtained by fluctuating along the least favorable submodel satisfies
$
P_n D^*_{f_n^1} = o_p(n^{-1/2}).
$
Consequently, $f_n^1$ also attains the same asymptotic efficiency, while preserving the convergence rate and sectional variation norm of the initial HAL-MLE. We formally state the theorem as follows.

\begin{theorem} [Asymptotic Efficiency of Single-step HAL‐TMLE] \label{thm:single-step-hal-tmle}
Suppose we impose a single targeting step to $f_n$ and yield $f_n^1$. And suppose that the second and third assumptions hold for $f_n^1$, that is,
\begin{enumerate}
\item (Negligible Remainder) \(R(P_{f_n^1},P_0)=o_p\bigl(n^{-1/2}\bigr)\).
\item (\(L^2(P_0)\)-continuity of EIC) \(P_0\{(D^*_{f_n^1}-D^*_{f_0})^2\}
      = o_p\bigl(1\bigr)\).
\end{enumerate}
Then $P_n D^*_{f_n^1} = o_p(n^{-1/2})$, and thus the single-step HAL-TMLE is asymptotically efficient.
\end{theorem}

Here are some examples for which we can compute this TMLE:

\begin{example}[Survival function and CDF]
For the marginal survival function $\Psi^F(f)=\int_{x_0}^1 p_f(y)dy$ at $x_0$, we have \(D^*_f(x) = I(x > x_0) - \mathbb{E}_{P_{f}}\bigl[I(x > x_0)\bigr].\) We can thus select $h = I(x > x_0)$. Note that, technically, we need $\mathbb{E}_{P_{f_n}}[h] = 0$, but the constant will cancel out by the normalizing constant. Similarly, for CDF $F_f$ at $x_0$ we can select $h=I(x\leq x_0)$.
\end{example}

\begin{example}[Median]
For the median $\Psi^F(f)=F_f^{-1}(0.5)$ we have 
\(D^*_f(x) = \frac{\frac12 - I(x < F_f^{-1}(0.5))}{f(F_f^{-1}(0.5)) },\) and thus $h = I(X< F_f^{-1}(0.5))$. We need to estimate the median based on the initial and then do the targeting. Note that the targeting of median is trivially generalized to any percentile. Since the density estimation changes from $f_n$ to $f_n^k$ after the $k$ steps targeting, we can iteratively do the targeting of the median. However, in Theorem~\ref{thm:single-step-hal-tmle}, we show that a single step targeting already guarantees $P_n D^*_{f_n^1} = o_p(n^{-1/2})$, and thus the asymptotic efficiency holds.
\end{example}

\begin{example}[$k$-th order moments]
For the $k$-th order moment $\Psi^F(f)=\int x^k p_f(x)\mathrm{d}x$ we have
\(D^*_f(x) = x^k - \Psi^F(f),\) and thus $h = x^k$. Note that this is just the parametric basis. This implies a motivation for not penalizing the parametric part, like LAS and the classic TF, in the sense that the resulting HAL fit would be targeting all its moments.
\end{example}
Note that the LLFPs for survival probability and moments are actually their universal least favorable paths (ULFP), and their corresponding single-step HAL-TMLE is the one-step HAL-TMLE in \citet{van2017generally, van2016one, van2018targeted}, and thus we have $P_n D^*_{f_n^1} = 0$ exactly. Also, by \citet{van2017generally}, $f_n^1$ and $f_n$ preserve the same convergence rate to $f_0$. Thus, we can extend the arguments of the assumptions of the plug-in HAL-MLE to the single-step HAL-TMLE. 

\subsection{Theoretical Property and Variance Estimation}
\label{subsec: EIC-based-variance-estimation}
As established in Section~\ref{Plug-in HAL-MLE} and Section~\ref{HAL-TMLE}, the score equation is solved at least at level $o_p(n^{-1/2})$. Thus, both estimators are asymptotically efficient and asymptotically linear, with their proofs unified in Section S.6 of the Supplementary Material \citep{hou2026supplement}. Since TMLE is asymptotically linear with influence curve $D^*_{f_n^1}$, variance estimation for $\Psi^F(f_n)$ can be obtained directly by computing the sample variance of the estimated $D^*_{f_n^1}(x_i)$.